\documentclass[10pt,a4paper]{amsart}
\usepackage{amsmath,amssymb,amsthm,graphicx}
% Retained preamble for the Stacks Project chapter divisors.tex (Divisors).
% The chapter file pinned-original-source.tex carries only \input{preamble} on its first line, so no
% preamble text lives inside the chapter itself.  The Stacks Project uses one single shared preamble for
% every chapter; the snapshot stages it verbatim in sources/stacks-chow/preamble.tex.  The block below is
% copied mechanically, line for line, from that file's lines 14-18 (the Xy-pic setup) and lines 154-217
% (the native theorem declarations together with the mathematical macros).  Nothing was added, removed or
% redefined.  The wrapper class and the amsmath/amssymb/amsthm/graphicx base packages are supplied by the
% assembler, so the class/\externaldocument part of the shared preamble is not repeated here.
% For commutative diagrams we use Xy-pic
\usepackage[all]{xy}

% We use 2cell for 2-commutative diagrams.
\xyoption{2cell}
% Theorem environments.
%
\theoremstyle{plain}
\newtheorem{theorem}[subsection]{Theorem}
\newtheorem{proposition}[subsection]{Proposition}
\newtheorem{lemma}[subsection]{Lemma}

\theoremstyle{definition}
\newtheorem{definition}[subsection]{Definition}
\newtheorem{example}[subsection]{Example}
\newtheorem{exercise}[subsection]{Exercise}
\newtheorem{situation}[subsection]{Situation}

\theoremstyle{remark}
\newtheorem{remark}[subsection]{Remark}
\newtheorem{remarks}[subsection]{Remarks}

\numberwithin{equation}{subsection}

% Macros
%
\def\lim{\mathop{\mathrm{lim}}\nolimits}
\def\colim{\mathop{\mathrm{colim}}\nolimits}
\def\Spec{\mathop{\mathrm{Spec}}}
\def\Hom{\mathop{\mathrm{Hom}}\nolimits}
\def\Ext{\mathop{\mathrm{Ext}}\nolimits}
\def\SheafHom{\mathop{\mathcal{H}\!\mathit{om}}\nolimits}
\def\SheafExt{\mathop{\mathcal{E}\!\mathit{xt}}\nolimits}
\def\Sch{\mathit{Sch}}
\def\Mor{\mathop{\mathrm{Mor}}\nolimits}
\def\Ob{\mathop{\mathrm{Ob}}\nolimits}
\def\Sh{\mathop{\mathit{Sh}}\nolimits}
\def\NL{\mathop{N\!L}\nolimits}
\def\CH{\mathop{\mathrm{CH}}\nolimits}
\def\proetale{{pro\text{-}\acute{e}tale}}
\def\etale{{\acute{e}tale}}
\def\QCoh{\mathit{QCoh}}
\def\Ker{\mathop{\mathrm{Ker}}}
\def\Im{\mathop{\mathrm{Im}}}
\def\Coker{\mathop{\mathrm{Coker}}}
\def\Coim{\mathop{\mathrm{Coim}}}

% Boxtimes
%
\DeclareMathSymbol{\boxtimes}{\mathbin}{AMSa}{"02}

%
% Macros for moduli stacks/spaces
%
\def\QCohstack{\mathcal{QC}\!\mathit{oh}}
\def\Cohstack{\mathcal{C}\!\mathit{oh}}
\def\Spacesstack{\mathcal{S}\!\mathit{paces}}
\def\Quotfunctor{\mathrm{Quot}}
\def\Hilbfunctor{\mathrm{Hilb}}
\def\Curvesstack{\mathcal{C}\!\mathit{urves}}
\def\Polarizedstack{\mathcal{P}\!\mathit{olarized}}
\def\Complexesstack{\mathcal{C}\!\mathit{omplexes}}
% \Pic is the operator that assigns to X its picard group, usage \Pic(X)
% \Picardstack_{X/B} denotes the Picard stack of X over B
% \Picardfunctor_{X/B} denotes the Picard functor of X over B
\def\Pic{\mathop{\mathrm{Pic}}\nolimits}
\def\Picardstack{\mathcal{P}\!\mathit{ic}}
\def\Picardfunctor{\mathrm{Pic}}
\def\Deformationcategory{\mathcal{D}\!\mathit{ef}}

% --- appended 2026-10-12 (Chief): the project own definitions of the reference/slogan/history
% display environments, copied mechanically from the shared Stacks preamble lines 8-12
% (\usepackage{verbatim} plus the three \newenvironment lines, which the chapter preamble keeps
% verbatim from the project preamble). Nothing added, removed or redefined: these environments
% carry no mathematical content and the project itself typesets them as comments. Without these
% three lines a counted statement whose own text uses them fails with Environment slogan undefined.
\usepackage{verbatim}
\newenvironment{reference}{\comment}{\endcomment}
\newenvironment{slogan}{\comment}{\endcomment}
\newenvironment{history}{\comment}{\endcomment}
\usepackage{cleveref}
\crefname{originaldefinition}{original definition}{original definitions}
\Crefname{originaldefinition}{Original definition}{Original definitions}
\makeatletter
\newlabel{definition-pullback-effective-Cartier-divisor}{{01WV}{source}{}{sourcepointer-1}{}}
\newlabel{definition-pullback-effective-Cartier-divisor@cref}{{[originaldefinition][0][]01WV}{[1][0][]source}{}{}{}}
\makeatother
\crefname{originaldefinition}{original definition}{original definitions}
\Crefname{originaldefinition}{Original definition}{Original definitions}
\makeatletter
\newlabel{definition-pullback-meromorphic-sections}{{02OT}{source}{}{sourcepointer-2}{}}
\newlabel{definition-pullback-meromorphic-sections@cref}{{[originaldefinition][0][]02OT}{[1][0][]source}{}{}{}}
\makeatother
\crefname{originallemma}{original lemma}{original lemmas}
\Crefname{originallemma}{Original lemma}{Original lemmas}
\makeatletter
\newlabel{lemma-blowing-up-affine}{{0804}{source}{}{sourcepointer-3}{}}
\newlabel{lemma-blowing-up-affine@cref}{{[originallemma][0][]0804}{[1][0][]source}{}{}{}}
\makeatother
\crefname{originallemma}{original lemma}{original lemmas}
\Crefname{originallemma}{Original lemma}{Original lemmas}
\makeatletter
\newlabel{lemma-characterize-effective-Cartier-divisor}{{01WS}{source}{}{sourcepointer-4}{}}
\newlabel{lemma-characterize-effective-Cartier-divisor@cref}{{[originallemma][0][]01WS}{[1][0][]source}{}{}{}}
\makeatother
\title{A blowup pulls back effective Cartier divisors and meromorphic functions}
\author{The Stacks Project Authors}
\date{stacks-project:divisors.tex:local-snapshot-20261012; GNU FDL 1.2 or later}
\begin{document}
\maketitle
\setcounter{subsection}{1}
\noindent\textit{Editorial note.} All delivered mathematical text is the author's own, copied from the pinned Stacks Project chapter file divisors.tex; no statement or proof step was rewritten, reordered or removed.\par\medskip
\noindent\textit{Editorial note.} Carrier window: statement and proof both lie inside the window 8000-10000 granted to this lane by CONTROL 8.83.\par\medskip
\noindent\textit{Editorial note.} The proof's algebra fact is proved in place, so the card is self-contained apart from the two reduction lemmas recorded as navigation pointers.\par\medskip
\noindent\textit{Editorial note.} Preamble: sources/stacks-divisors/retained-preamble.tex (88 lines, verbatim present).\par\medskip
\noindent\textit{Allowed standard prerequisites (editorial).}\begin{itemize}
\item Local description of a blowup by affine blowup algebras. Lemma lemma-blowing-up-affine, used at proof lines 8577-8579. Reduces the assertion to the affine case.
\item Local characterisation of effective Cartier divisors. Lemma lemma-characterize-effective-Cartier-divisor, used at proof lines 8577-8579. Turns the pullback statement into the nonzerodivisor statement.
\item Definitions of the pullback of an effective Cartier divisor and of meromorphic sections. Definitions definition-pullback-effective-Cartier-divisor and definition-pullback-meromorphic-sections, used at proof lines 8570-8574. Fix what has to be shown defined.
\end{itemize}\medskip
\ifdefined\hypertarget\hypertarget{sourcepointer-1}{}\fi
\noindent\textit{Source reference.} Original definition 01WV: focused-seeds/AGP/sources/stacks-snapshot-20261011/tags.txt, source lines 2496--2496. Navigation pointer into the Stacks Project; the standard fact it supplies is applied inside the counted proof span and the argument resting on it stays there.\par\medskip
\ifdefined\hypertarget\hypertarget{sourcepointer-2}{}\fi
\noindent\textit{Source reference.} Original definition 02OT: focused-seeds/AGP/sources/stacks-snapshot-20261011/tags.txt, source lines 3528--3528. Navigation pointer into the Stacks Project; the standard fact it supplies is applied inside the counted proof span and the argument resting on it stays there.\par\medskip
\ifdefined\hypertarget\hypertarget{sourcepointer-3}{}\fi
\noindent\textit{Source reference.} Original lemma 0804: focused-seeds/AGP/sources/stacks-snapshot-20261011/tags.txt, source lines 10169--10169. Navigation pointer into the Stacks Project; the standard fact it supplies is applied inside the counted proof span and the argument resting on it stays there.\par\medskip
\ifdefined\hypertarget\hypertarget{sourcepointer-4}{}\fi
\noindent\textit{Source reference.} Original lemma 01WS: focused-seeds/AGP/sources/stacks-snapshot-20261011/tags.txt, source lines 2493--2493. Navigation pointer into the Stacks Project; the standard fact it supplies is applied inside the counted proof span and the argument resting on it stays there.\par\medskip
\begin{lemma}
\label{lemma-blow-up-pullback-effective-Cartier}
Let $X$ be a scheme. Let $b : X' \to X$ be a blowup of $X$ in a closed
subscheme. The pullback $b^{-1}D$ is defined
for all effective Cartier divisors $D \subset X$
and pullbacks of meromorphic functions are defined for $b$
(Definitions
\ref{definition-pullback-effective-Cartier-divisor} and
\ref{definition-pullback-meromorphic-sections}).
\end{lemma}

\begin{proof}
By Lemmas \ref{lemma-blowing-up-affine} and
\ref{lemma-characterize-effective-Cartier-divisor}
this reduces to the following algebra fact:
Let $A$ be a ring, $I \subset A$ an ideal, $a \in I$, and $x \in A$
a nonzerodivisor. Then the image of $x$ in $A[\frac{I}{a}]$ is a
nonzerodivisor. Namely, suppose that $x (y/a^n) = 0$ in $A[\frac{I}{a}]$.
Then $a^mxy = 0$ in $A$ for some $m$. Hence $a^my = 0$ as $x$ is a
nonzerodivisor. Whence $y/a^n$ is zero in $A[\frac{I}{a}]$ as desired.
\end{proof}

\end{document}
